\documentclass[11pt]{article}
\usepackage[a4paper,margin=0.8in]{geometry}
\usepackage{amsmath,amssymb,booktabs,graphicx}
\usepackage[colorlinks=true,linkcolor=blue,urlcolor=blue]{hyperref}
\setlength{\emergencystretch}{2em}
\graphicspath{{../output/moving_comparison/}}
\title{MovingThreeFluidSim: implementation checks and sample comparisons}
\author{Implementation on the upstream 49dd74e baseline}
\date{22 September 2026}
\begin{document}
\maketitle

\section{Outcome and comparison convention}
\emph{Historical centred-flux checkpoint. The updated numerical-scheme note
documents the alternating-flux correction and new runs that remove the
checkerboard failure. The results below are retained for comparison, not
as the current implementation's validation status.}

The implementation checkpoint is \texttt{17285f4}, on branch
\texttt{derivation/moving-three-fluid}.
The optimized $12N$ solver is implemented, with scalar half-bandwidths
$(13,14)$ and one banded linear equation per step. The companion scheme
note and Mathematica compiler include the implementation corrections.
\textbf{The long-run validation gate is not yet met.} The single fluid
and its equal-mass split agree with one another, but the unequal-mass
comparison differs substantially and very dilute components develop
oscillations that make the adaptive timestep impractically small.
No density floor, timestep retry, deleted trace component or shifted
collapse time was used to conceal these failures.

Both solvers start from identical density, internal energy, grid and
particle masses, with identical $c_1,c_2,c_4$. All tests here have $q=0$,
no formation and no imposed stripping. Separate unit tests exercise $q>0$.
The canonical case is the no-formation baseline of \texttt{main.cpp};
the 1000-cell case is its resolution test. The Yiming cases retain their
150-cell initialization. Their stored masses imply a different particle
count from the canonical $10^6$-star initialization; we do not silently
change them. The inertia coefficient uses
\begin{equation}
 \epsilon=(3m_s\ln\Lambda_{sd})^2,\qquad
 \ln\Lambda_{sd}=\ln\left(\frac{0.8M_{\rm tot}}{m_s+m_d}\right).
\end{equation}
It is $1.2533\times10^{-9}$ for the single/split cases,
$1.3835\times10^{-9}$ for the AB-like baseline and
$1.8475\times10^{-10}$ for the canonical case.

Plots compare actual accepted code times and common cell-centre radii.
The plotted one-dimensional dispersion is $\sigma/\sqrt{u_0}=\sqrt{2u/3}$.
Neither time nor density is rescaled to align collapse. Solid/dashed
profile curves denote hydrostatic/moving solvers; rainbow colour denotes
time. Both solvers may overshoot the final time, and temporal interpolation
is confined to their common output interval. The moving solver has open
material/thermal boundaries, unlike the hydrostatic outer closure. Its
stellar mass loss is below $7\times10^{-6}$ in these tests, so a large
core discrepancy cannot be explained by a large global escaped mass.

\section{Checks that pass}
The 34 Mathematica checks pass, including the optimized storage permutation,
all boundary regimes, nonlinear finite-difference Jacobians and the assembled
continuity ledger. Every C++ matrix entry and RHS entry was compared with
the independent symbolic derivation in a five-cell fixture containing
fan, supersonic and vacuum outer branches. The maximum scaled discrepancies
were $3.73\times10^{-16}$ and $3.39\times10^{-16}$, respectively.

The C++ checks cover band versus dense solution, central regularity,
interior equilibrium residuals, the tidal source and outward evolution,
single/split equivalence, reference depletion, local thermal agreement with
the existing conduction operator, invalid input and parameter serialization.
The existing hydrostatic, Statler and sample checks and all 12 Python tests
pass. These short checks do not establish long-run stability.

Across the reported runs, relative mass-ledger errors remain below
$3\times10^{-15}$. The per-step gas-energy/work residual is below
$5\times10^{-17}$. The latter compares random-plus-bulk energy with the
included gravitational/reference work, heating and exported energy;
it does not assert exact conservation of total gravitational energy.
For 50, 100, 200, 400 and 800 cells, measured step times were approximately
0.47, 0.93, 1.88, 3.90 and 8.42 ms in the short benchmark. Band storage
doubles exactly with $N$, consistent with the explicit $O(N)$ loops and
fixed-band factorization.

\section{Evolution comparisons}
Table~\ref{tab:results} gives central stellar ratios at the last common
time. A partial run is not a successful integration to $t=5$.
\begin{table}[ht]
\centering
\begin{tabular}{lrrrr}
\toprule
Case & $N$ & Last common $t$ & $\rho_{s,c}^{\rm mov}/\rho_{s,c}^{\rm hyd}$
& $\sigma_{s,c}^{\rm mov}/\sigma_{s,c}^{\rm hyd}$\\
\midrule
Single fluid &150&5.0001&0.5072&0.9408\\
Equal-mass split &150&5.0001&0.5072&0.9408\\
AB-like, no stripping &150&2.8002&0.1580&0.9021\\
Canonical, partial &500&3.1371&0.9755&1.0000\\
Canonical refined, partial &1000&2.7291&0.9932&1.0005\\
\bottomrule
\end{tabular}
\caption{Ratios use unshifted code time. Both canonical runs requested $t=5$
but stopped at the 12000-step safety limit after trace-fluid oscillations
made their timestep small.}\label{tab:results}
\end{table}

For the 150-cell single fluid, central density agrees within 10\% for
about 62\% of the interval $0<t<5$; dispersion agrees within 10\% throughout.
The late contraction is slower in the moving calculation. The equal split
does not introduce a new effect: its total central density differs from
the unsplit moving calculation by less than $10^{-9}$, and mass-weighted
dispersion by less than $8\times10^{-11}$ over the sampled common history.
Thus splitting a physically identical fluid passes the consistency check.

\begin{figure}[ht]
\centering
\IfFileExists{../output/moving_comparison/matched_single/central_evolution.pdf}{%
\includegraphics[width=.9\textwidth]{matched_single/central_evolution.pdf}}{%
\fbox{\parbox{.85\textwidth}{Generate the matched-single comparison to insert this figure.}}}
\caption{Single-fluid evolution at the original 150-cell resolution.
Agreement at early times does not imply agreement close to collapse.}
\end{figure}

The canonical stellar component tracks the hydrostatic result much more
closely before the trace-fluid failure. Refinement improves that stellar
agreement, but does not cure the failure: the 1000-cell run stalls earlier
at the same bounded step budget. In the 500-cell run, strong cell-to-cell
structure develops in the extremely dilute DM component. Its density and
temperature have anticorrelated oscillations, accompanied by alternating
radial velocities. The timestep controller is responding to that state;
the million-step diagnostic did not recover useful progress. Small linear
residuals and accurate mass budgets do not validate an oscillatory solution.

\begin{figure}[ht]
\centering
\IfFileExists{../output/moving_comparison/review_canonical/all_components_final.pdf}{%
\includegraphics[width=\textwidth]{review_canonical/all_components_final.pdf}}{%
\fbox{\parbox{.9\textwidth}{Generate the canonical comparison to insert this figure.}}}
\caption{All three components at the end of the partial 500-cell run.
The mass-dominant stellar agreement alone would conceal the DM failure.
The bottom row shows raw moving-fluid cell velocities, not a velocity
inferred from the hydrostatic calculation.}
\end{figure}

For the AB-like unequal-mass case, the moving calculation segregates less
strongly. At $t=2.8$, its central DM density is 2.89 times the hydrostatic
value, while its central DM dispersion is 0.746 times that value.
The requested majority-of-evolution agreement is therefore not obtained
for this case. This cannot be presented as a validated reproduction.

\section{The hydrostatic reference is not a conservative ground truth}
The upstream realignment logarithmically interpolates density, recomputes
enclosed masses, then rebuilds pressure by a hydrostatic integration and
sets $u=3P/(2\rho)$. It does not conservatively transport mass and entropy
through that remap. We left this existing implementation unchanged.

An operator audit of the initial AB-like case isolates $c_1$ exchange by
setting $c_2=0$. With $\Delta t=10^{-3}$, realignment changes total mass by
$-2.17\times10^{-5}$ and random energy by $-2.30\times10^{-6}$, relative to
the immediately preceding relaxed state. These are absolute code-unit
changes, not full-step physical energy losses.

There is also a finite-grid, zero-step effect. At $\Delta t=0$, the old
conduction boundary still sets $u_0=u_1$, changing central $u_s$ by
$-9.6372\times10^{-6}$ in this initial profile. Subsequent relaxation and
realignment change the profile despite zero elapsed time. In this audit
the realignment mass and random-energy changes are
$-3.4186\times10^{-7}$ and $-3.5137\times10^{-8}$. Therefore dividing these
initial-step changes by an ever smaller $\Delta t$ does not measure a
convergent physical rate. The moving finite-volume centre instead retains
the first cell's energy equation and imposes zero face flux.

In addition, the per-fluid hydrostatic projection uses sums of masses at
the same shell index while fluid radii may differ. For relative shell
displacements $\xi_f$, a common-radius treatment contains
\begin{equation}
 \delta M_{\rm tot}(R_f)=R_f^2\sum_g\rho_g(\xi_f-\xi_g).
\end{equation}
This vanishes for an identical co-moving split but not for unequal flows.
Its omission in a per-fluid Jacobian, followed by a remap and pressure
reset, is another distinction to audit. The present tests do not isolate
its quantitative share of the AB discrepancy. Nor do they establish that
the new scheme is correct simply because the reference has deficiencies.

\section{Remaining acceptance work}
The implementation should remain experimental. The next numerical task
is an entropy-compatible, low-Mach, well-balanced interface treatment
that remains admissible when a trace component departs strongly from the
fixed reference. A smaller timestep, deleting that component or adding a
temperature floor is not a demonstrated cure. Empty-cell transitions also
remain outside the positive-state linearization.

Separately, a conservative hydrostatic remap and compatible centre treatment
are needed before demanding close multi-fluid agreement with that reference.
These are changes to the existing solver and were not made in this task.
Reproduction commands, binary output conventions and the explicit current
limitations are documented with the C++ implementation. The generated PDFs
include density/dispersion profiles, central histories and all-component
diagnostics; partial histories are labelled as such.
\end{document}
